\subsection{An actual cubic invariant constructed using a prime in a residue class}
\label{sec:hessian-fermat-prime}

\paragraph{P-F49 (Hessian--Cayley nonvanishing, with exact prime bounds).}
Let $n\ge2$, let $p\equiv1\pmod n$ be prime, and put
$D=p-1$, $h=D/n$. For a cubic $F$ in $n$ variables set
\[
       H_F(x)=6^{-n}\det\left(
                   \frac{\partial^2F}{\partial x_i\partial x_j}
                                      \right)_{1\le i,j\le n}.
                                                        \tag{P-F49a}
\]
There is an explicit polynomial $I_{n,p}$ in the cubic
coefficients, defined over $\mathbb Q$, such that
\[
 \deg I_{n,p}=n(p-1),\qquad
 I_{n,p}(gF)=(\det g)^{3(p-1)}I_{n,p}(F),\qquad
 I_{n,p}(F_n)\ne0,\quad F_n=\sum_{i=1}^n x_i^3.
                                                        \tag{P-F49b}
\]
It is obtained by applying Cayley's degree-$n$ determinant
contraction on forms of degree $D$ to $H_F^h$.
In the actual cubic plethysm notation, $g_n=\gcd(n,3)$,
this proves
\[
              f_n(g_n(p-1))>0,
           \qquad \theta_{3(p-1)}^{\wedge n}\ne0.
                                                        \tag{P-F49c}
\]
Unconditionally one can choose $p\le Cn^{5.18}$ with
an effective absolute constant $C$. For $n\ge4$, the
Dirichlet-GRH least-prime theorem gives the sharper choice
\[
      p\le(\varphi(n)\log n)^2,
      \quad g_n(p-1)<g_n(\varphi(n)\log n)^2,
      \quad\deg I_{n,p}<n^3(\log n)^2.                 \tag{P-F49d}
\]
The supplied Dirichlet statements A-DIR-NV/A-DIR-Z
discharge exactly this GRH premise in the present work.
These are actual point-specific invariant occurrences;
they do not assert a polynomial-size evaluation circuit,
an invariant basis, or a bound on the strict PH2/SH indices.

\paragraph{Cayley's contraction and its normalization.}
For a degree-$D$ form
$G(x)=\sum_{|\alpha|=D}g_\alpha x^\alpha$, let its
symmetric tensor entries be
\[
       G_{i_1\cdots i_D}=g_\alpha\frac{\alpha!}{D!},
       \qquad \alpha!=\prod_{a=1}^n\alpha_a!,
                                                        \tag{P-F49e}
\]
when the index list has multiplicity vector $\alpha$.
Define
\[
 \mathcal C_{D,n}(G)=
 \sum_{\sigma_1,\ldots,\sigma_D\in\mathfrak S_n}
 \left(\prod_{b=1}^D\operatorname{sgn}(\sigma_b)\right)
 \prod_{a=1}^n G_{\sigma_1(a),\ldots,\sigma_D(a)}.
                                                        \tag{P-F49f}
\]
There are $n$ tensor factors in each summand, so this
is a degree-$n$ polynomial in the coefficients of $G$.
Contracting each column of $n$ indices with a determinant
gives
\[
             \mathcal C_{D,n}(gG)
                          =(\det g)^D\mathcal C_{D,n}(G).
                                                        \tag{P-F49g}
\]
This is precisely the invariant in B\"urgisser--Ikenmeyer,
equation~(3.8), printed p.~11, with its transformation
law in Theorem~3.18(1), p.~14. We retain the formula to
make the normalization used below explicit.

\paragraph{The determinant coefficient in prime characteristic.}
For an $n\times n$ matrix $X=(X_{ij})$ of indeterminates,
write $c_L$ for the coefficient of $X^L$ in
$(\det X)^{p-1}$, where every row and column of $L$ sums
to $p-1$. Then
\[
       c_L\equiv\frac{((p-1)!)^n}{\prod_{i,j}L_{ij}!}
                                 \not\equiv0\pmod p.  \tag{P-F49h}
\]
This is Glynn's determinant coefficient identity, stated
with the congruence in Itoh--Shimoyoshi,
Theorem~1.1 and the following Remark, printed p.~1,
\url{https://arxiv.org/pdf/2104.00838}.
Here is the short form of its proof needed for the construction.

Work over the algebraic closure of $\mathbb F_p$ in
\[
        A=\overline{\mathbb F}_p[z_1,\ldots,z_n]
                                      /(z_1^p,\ldots,z_n^p).
\]
Every linear substitution preserves this ideal by
the characteristic-$p$ power identity. The top degree
$n(p-1)$ is a line spanned by $\prod_i z_i^{p-1}$.
Diagonal substitutions act on it by $(\det X)^{p-1}$.
An elementary substitution $z_i\mapsto z_i+a z_j$
acts trivially: every extra term has exponent at least
$p$ in $z_j$ and vanishes. Diagonal and elementary
matrices generate $GL_n$, so the top-line character is
$(\det X)^{p-1}$ for every invertible matrix. Since the
field is infinite this is an identity of polynomials
in the matrix entries. Therefore
\[
 [z_1^{p-1}\cdots z_n^{p-1}]
     \prod_{j=1}^n\left(\sum_i X_{ij}z_i\right)^{p-1}
                          =(\det X)^{p-1}.             \tag{P-F49i}
\]
Expanding the left side proves (P-F49h). Every entry
of $L$ is between zero and $p-1$, so all the factorials
in its denominator are units modulo $p$.

\paragraph{Nonzero evaluation and the exact cubic character.}
Put $G_0=(x_1\cdots x_n)^h$. Its tensor entries vanish
unless every symbol occurs $h$ times, and in that case
equal $(h!)^n/D!$. Consequently (P-F49f) gives the exact
identity over $\mathbb Q$
\[
 \mathcal C_{D,n}(G_0)
       =\left(\frac{(h!)^n}{D!}\right)^n
              [\prod_{i,j}X_{ij}^{h}](\det X)^D.
                                                        \tag{P-F49j}
\]
The coefficient on the right is nonzero modulo $p$
by (P-F49h). Since $h,D<p$, all displayed factorials
are units there; in fact (P-F49j) is congruent to one
modulo $p$. In particular it is nonzero in characteristic
zero. This argument uses the reduction of the specified
monomial evaluation; it does not require reducing the
normalizing factor $6^{-n}$ when $p=3$.

Define $I_{n,p}(F)=\mathcal C_{D,n}(H_F^h)$.
The Hessian determinant has degree $n$ in the cubic
coefficients and degree $n$ in $x$, so $H_F^h$ has
coefficient degree $nh=D$ and variable degree $D$.
The coefficient degree of $I_{n,p}$ is therefore $nD$.
For the convention $(gF)(x)=F(g^t x)$, the chain rule
gives
\[
            H_{gF}(x)=(\det g)^2H_F(g^t x).
                                                        \tag{P-F49k}
\]
After taking the $h$th power and applying (P-F49g),
the determinant exponent is $2hn+D=3D$. This proves
the transformation law in (P-F49b). Finally
$H_{F_n}=x_1\cdots x_n$, so (P-F49j) proves its nonzero
Fermat value. Clearing rational denominators gives an
integral-coefficient invariant if desired, without
changing its degree or its characteristic-zero value.

The resulting invariant is a nonzero element of the
degree-$nD$ invariant space. Dualizing the finite-dimensional
representation preserves the dimension of its invariant
space, and $nD=ns/g_n$ for $s=g_nD$. This proves the
first assertion of (P-F49c); P-F47 proves the second.

\paragraph{The exact analytic premise and the degree bounds.}
Xylouris, \emph{On the least prime in an arithmetic
progression and estimates for the zeros of Dirichlet
$L$-functions}, Theorem~1.1, printed p.~66, proves
$P(a,m)\le C m^{5.18}$ for reduced residue classes,
with an effectively computable absolute constant $C$.
Applying it to $a=1,m=n$ gives the unconditional
polynomial choice in (P-F49b)--(P-F49c).

Bach--Sorenson, \emph{Explicit bounds for primes in
residue classes}, Theorem~5.3, pp.~1731--1732, gives
$P(a,m)<2(m\log m)^2$ under its ERH premise. Its
arithmetic-progression specialization takes
$E=\mathbb Q(\zeta_m)$; the relevant zeta function is
the product of primitive Dirichlet $L$-functions, as
explained in their Section~2, pp.~1719--1720.
Thus the supplied Dirichlet theorem suffices for this
particular specialization; arbitrary number-field
GRH is not being imported.

We use the sharper bound of Lamzouri--Li--Soundararajan,
\emph{Conditional bounds for the least quadratic
non-residue and related problems}, Corollary~1.2,
printed p.~3, which states
$P(a,m)\le(\varphi(m)\log m)^2$ for $m>3$.
Their Section~2, pp.~4--5, uses RH for $\zeta$ and
GRH for primitive Dirichlet $L$-functions. These lie
within the supplied A-DIR-NV/A-DIR-Z scope. Setting
$a=1,m=n\ge4$ gives (P-F49d). For the two smaller
dimensions one may take $p=3$ for $n=2$, and $p=7$
for $n=3$, obtaining stretches $2$ and $18$ respectively.
No optimality in those small cases is asserted.

The primary URLs for the analytic statements are
\begin{itemize}
\item Xylouris, Theorem~1.1, p.~66:
\url{https://www.impan.pl/shop/publication/transaction/download/product/83078}.
\item Bach--Sorenson, Theorem~5.3, pp.~1731--1732:
\url{https://doi.org/10.1090/S0025-5718-96-00763-6}.
\item Lamzouri--Li--Soundararajan, Corollary~1.2,
p.~3, and Section~2, pp.~4--5:
\url{https://arxiv.org/pdf/1309.3595}.
\end{itemize}

\paragraph{What has and has not become constructive.}
The formulas specify a genuine invariant and a proved
nonzero evaluation for every $n$, with a polynomial
degree bound even without the supplied theorem and a
sharper explicit bound using that theorem. The prime
may be found by searching the bounded progression.
However, the displayed Cayley contraction has
$(n!)^{p-1}$ summands before simplification. Its polynomial
degree does not establish a polynomial arithmetic-circuit
size or a positive counting model for multiplicities.
The theorem closes the polynomial-degree Fermat
nonvanishing question posed in P-F47, not the distinct
strict-constituent PH2 or SH question.
